% AUTO-GENERATED. DO NOT EDIT.
% Generator: hpc/paper_artifacts/build_tab_baselines.py
% Benchmark: EvEMTBench v1.1.0
% Inputs: (none — static specification)
\begin{table*}[t]
\centering
\small
\caption{Reference baselines: a trivial floor, a conventional-protection anchor,
classical \ac{ml}, and deep learning. The \ac{gnn} is planned for the cross-topology
protocol.}
\label{tab:baselines}
\begin{tabular}{L{2.2cm} L{2.8cm} L{3.0cm} L{4.6cm}}
\toprule
\textbf{Baseline} & \textbf{Family} & \textbf{Input} & \textbf{Notes} \\
\midrule
majority & trivial floor & labels only & per-task chance level \\
threshold & conventional (overcurrent) & local current & binary \ac{fd}, \emph{local} only; the conventional-protection anchor \\
\ac{rf} & classical \ac{ml} & per-channel summary features & tree-based ensemble \\
\ac{mlp} & deep (flattened window) & raw window & no hand-crafted features \\
\ac{gru} & deep (sequence) & raw window & sequence model \\
\ac{cnn} & deep (1-D waveform) & raw window & waveform model \\
\ac{resnet} & deep (1-D residual) & raw window & residual waveform model \\
\ac{gnn} \planned & deep (topology-aware) & grid graph & cross-topology contender \\
\bottomrule
\end{tabular}
\end{table*}
